\chapter[Die Kraft des Wortes Gottes]{Die Kraft des Wortes Gottes}
\quellseite{10}
Es war um das Jahr 1665, daß der Prediger Johann Flavel, der in der englischen Handelsstadt Dartmouth durch sein gewaltiges evangelisches Zeugniß viele Seelen aus dem geistlichen Tod zum Leben in Christo weckte, an einem Sonntag über die Worte Pauli 1 Cor. 16, 22. predigte: „So jemand den Herrn Jesum Christum nicht lieb hat, der sei Anathema, Maharam Motha“, das heißt: der sei verflucht, oder zum Tode verbannt, wenn der Herr zum Gericht über die Gottlosen kommt. Die Predigt war unaussprechlich ernst und ergreifend und machte einen erschütternden Eindruck auf die Versammlung. Ein Herr aus hohem Stande wurde von seinen Gefühlen so übernommen, daß er ohnmächtig zu Boden fiel.

Unter den Anwesenden befand sich auch Lucas Short, damals ein junger Bursche von fünfzehn Jahren, der Sohn ehrbarer Bürgerleute aus Dartmouth, aber leichtsinnig, gedankenlos und weltlustig. Nicht lange nach dieser Predigt, die zwar augenblicklich einen Eindruck auf ihn machte, aber doch bald wieder vergessen zu sein schien, ließ er sich auf einem Schiff als Matrose anwerben und trieb sich eine Zeitlang auf dem Meere umher, bis er endlich, des Seelebens müde, in America ans Land stieg, sich da ein Grundstück erwarb und nun den Rest seiner Tage in der neuen Heimath verlebte.

Äußerlich ging es ihm gut. Er fand eine brave Lebensgefährtin, eine Anzahl gesunder und kräftiger Kinder wuchs ihm auf; sein Landgut war im blühenden Zustand; es schien zu seinem Glück kaum noch etwas zu fehlen. Ueber das alles erreichte Short ein Alter, das weit über die gewöhnliche Länge der menschlichen Lebenstage hinausging. Als er bereits hundert Jahre alt war, arbeitete er noch mit ungebrochener Kraft auf seinen Feldern; auch seine Geisteskräfte waren noch so frisch und ungeschwächt, als wäre er erst ein Fünfziger.

Aber ach, er war in Gefahr, zu denen gerechnet zu werden, die selbst im Angesicht des nahenden Todes in der Unbußfertigkeit verharren. Er ging in keine Kirche, las kein Wort Gottes, hielt mit den Seinen keine Morgen- und Abendandacht, und in seinem Hause hörte man nie die Stimme des Gebets. Wohl ging es äußerlich ehrbar zu. Da war kein Fluchen, kein Zank, kein Streit, kein Lärm des Weinsäufers oder dergleichen; aber die Seelen waren eben „ohne Christum, fremde und außer der Bürgerschaft Israel, und fremde von den Testamenten der Verheißung“ und daher „ohne Hoffnung und ohne Gott in der Welt“; und es bleibt in Ewigkeit wahr, was der Prophet sagt: „Die Gottlosen haben keinen Frieden.“

Eines Abends, als das Tagewerk vollendet war — es war im Jahr 1750 — saß der alte Short auf einem umgestürzten Baumstamm auf dem Feld und versank in tiefes Nachdenken, und wie es bei alten Leuten geht, verloren sich seine Gedanken weit zurück in die frühen Jugendtage und verweilten dabei sehr lange. Seiner längst verstorbenen Eltern Haus in Dartmouth, die Straßen und der Hafen der Stadt mit den ab und zu gehenden Schiffen, die Spiele und Tummelplätze seiner Jugend, die Schule, in der er mit seinen Altersgenossen gelernt hatte — alles, alles trat ihm mit ungewöhnlicher Lebendigkeit vor die Seele. Ein Name um den andern von denen, die er gekannt und geliebt hatte, drängte sich in sein Gedächtniß; wo waren sie jetzt alle? Er ging im Geist auf den Kirchhof und sah dort einen Grabhügel um den andern sich erheben, unter denen nun wohl alle, die er einst gekannt hatte, schon längst schliefen. Ueber den Kirchhof ragte die Kirche seiner Vaterstadt empor. Dort war er getauft und confirmirt; dort hatte er zum erstenmal unter tiefen Bewegungen des Herzens das heilige Abendmahl genossen. Und seitdem? Ach, seitdem hatte er sich nie wieder zum Tisch des Herrn genähert! Da erscheint vor seinem Geist der gewaltige Flavel, der Prediger, den er in den Tagen seiner Kindheit so oft gehört hat; es ist ihm, als sähe er ihn vor sich stehen mit der ehrwürdigen Gestalt, mit der aufgehobenen Hand, mit dem hellen, durchdringenden Auge. Er hört seine bewegliche Stimme! Und plötzlich tritt ihm jener Sonntag vor Augen, da er ihn fast zum letztenmal gehört hat, jener Sonntag, an welchem Flavel über den Text predigte: „So jemand den Herrn Jesum Christum nicht lieb hat, der sei Anathema, Maharam Motha!“ Der alte Short versinkt tiefer und tiefer in die Erinnerungen jenes unvergesslichen Tages, er hört noch einmal jene Predigt, er sieht die tiefbewegte Versammlung, sieht den vornehmen Herrn ohnmächtig hinsinken! Der Alte fühlt sich selbst einer Ohnmacht nahe. Zitternd erhebt er sich von dem Baumstamm und geht, er weiß selbst nicht wohin. Die Sonne neigt sich zum Untergang und mahnt ihn an das nahe Ende seines eignen Lebenstages. „So jemand den Herrn Jesum nicht lieb hat, der ist verflucht und zum Tod verbannt!“ Das Wort wird er nicht wieder los. Siehe, der Same dieses Wortes ist noch lebendig, obschon es nun schon 85 Jahre her ist, daß er es gehört hat.

Die Stunde Gottes war gekommen. Ein Thau vom Himmel und ein belebender Sonnenstrahl der Gnade fiel in sein Herz, und es dauerte nicht lange, da siehst du in der Kammer des alten Short die Bibel offen auf dem Tisch liegen; du hörst die Stimme des Rufens um Gnade und Sündenvergebung und dann auch die Stimme des Lobens und Dankens für das lebendig erfahrene Heil in Christo Jesu. Noch sechzehn Jahre lebte Lucas Short als ein Denkmal der Gnade Gottes und als ein Zeugniß der unvergänglichen Lebenskraft des göttlichen Wortes. Als ein Jünger Jesu ist er im Alter von 116 Jahren zur ewigen Ruhe eingegangen.

(Blüthen aus dem Garten Gottes.)
